% Portada de J2.3: una señal con un pico aislado, un tramo plano y un salto de nivel.
\documentclass[tikz,border=0pt]{standalone}
\usepackage{resonancia}
\begin{document}
\begin{tikzpicture}
  \clip[rounded corners=6pt] (0,0) rectangle (12,3.6);
  \shade[top color=rNaranja!14, bottom color=white] (0,0) rectangle (12,3.6);
  \draw[eje] (0.5,1.9) -- (11.6,1.9);
  % tramo 1 con pico
  \draw[medida, grueso] plot[smooth, domain=0.6:3.6, samples=120] (\x, {1.9 + 0.35*sin(800*(\x-0.6)) + 0.12*sin(1400*(\x-0.6))});
  \draw[color=rNaranja, line width=1.4pt] (2.1,1.9) -- (2.1,3.3);
  \fill[rNaranja] (2.1,3.3) circle (2.5pt);
  \node[etiqueta, color=rNaranja] at (2.1,3.5) {\small pico};
  % tramo plano
  \draw[medida, grueso] plot[smooth, domain=3.6:4.8, samples=60] (\x, {1.9 + 0.35*sin(800*(\x-0.6))*(4.8-\x)/1.2});
  \draw[color=rNaranja, line width=1.6pt] (4.8,1.9) -- (6.3,1.9);
  \node[etiqueta, color=rNaranja] at (5.55,1.45) {\small plano};
  % salto
  \draw[medida, grueso] plot[smooth, domain=6.3:7.6, samples=60] (\x, {1.9 + 0.3*sin(800*(\x-0.6))*(\x-6.3)/1.3});
  \draw[medida, grueso] plot[smooth, domain=7.6:11.4, samples=160] (\x, {2.55 + 0.35*sin(800*(\x-0.6)) + 0.12*sin(1400*(\x-0.6))});
  \draw[color=rNaranja, line width=1.4pt, dash pattern=on 3pt off 2pt] (7.6,1.9) -- (7.6,2.55);
  \node[etiqueta, color=rNaranja] at (8.6,1.25) {\small salto de nivel};
  \node[etiqueta, color=rGris] at (6.0,0.5) {\small detectar, decidir y dejar constancia};
\end{tikzpicture}
\end{document}
